\documentclass[11pt,a4paper]{article}
\usepackage[utf8]{inputenc}
\usepackage[T1]{fontenc}
\usepackage[polish]{babel}
\usepackage{geometry}
\geometry{margin=1.5cm}
\usepackage{tabularx}
\usepackage{longtable}

\begin{document}

\begin{flushright}
Zał. nr 2a
\end{flushright}

\noindent
Akademia Nauk Stosowanych w Elblągu \\
Instytut Informatyki Stosowanej im. Krzysztofa Brzeskiego \\
Kierunek studiów: Informatyka

\vspace{1em}

\noindent Student/ka: \dotfill \\
Nr albumu: \dotfill \\
Kierunek studiów: Informatyka \hfill Specjalność: \dotfill \\
Miejsce praktyki (instytucja): \dotfill \\
\dotfill \\
Termin realizacji praktyki: od \dots\dots\dots\dots\dots\ do \dots\dots\dots\dots\dots\ \hfill Liczba dni roboczych: 120

\vspace{2em}
\begin{center}
\textbf{PROGRAM PRAKTYKI ZAWODOWEJ}
\end{center}

\renewcommand{\arraystretch}{1.5}
\begin{longtable}{|c|p{8cm}|p{7cm}|}
\hline
\textbf{Lp.} & \textbf{Efekty kształcenia} & \textbf{Dział (komórka) / przykładowe prace wykonywane przez praktykanta} \\ \hline
\endfirsthead
\hline
\textbf{Lp.} & \textbf{Efekty kształcenia} & \textbf{Dział (komórka) / przykładowe prace wykonywane przez praktykanta} \\ \hline
\endhead
01 & Ma wiedzę na temat sposobu realizacji zadań inżynierskich dotyczących informatyki z zachowaniem standardów i norm technicznych & \\ \hline
02 & Zna technologie, narzędzia, metody, techniki oraz sprzęt stosowane w informatyce & \\ \hline
03 & Zna ekonomiczne, prawne skutki własnych działań podejmowanych w ramach praktyki oraz ograniczenia wynikające z prawa autorskiego i kodeksu pracy & \\ \hline
04 & Zna zasady bezpieczeństwa pracy i ergonomii w zawodzie informatyka & \\ \hline
05 & Pozyskuje informacje odnośnie technologii, metod, technik, sprzętu wymaganego do realizacji powierzonego zadania, posługując się rozmaitymi źródłami literaturowymi i zasobami publikowanymi w języku polskim jak i angielskim & \\ \hline
06 & W oparciu o kontakty ze środowiskiem inżynierskim zakładu, potrafi podnieść swoje kompetencje, wiedzę i umiejętności, co najmniej z dwóch zakresów: zadania dotyczące sprzętu i oprogramowania: np.: programowania, administrowanie siecią komputerową, konserwacja sprzętu i oprogramowania, bieżące usuwanie usterek, administrowanie zasobami informatycznymi, zakładu pracy / instytucji, (e)‑usługami. & \\ \hline
07 & Opracowuje dokumentację dotyczącą realizacji podejmowanych zadań w ramach praktyki, a także referuje ustnie prezentowane w niej zagadnienia & \\ \hline
08 & Potrafi zidentyfikować problem informatyczny występujący w zakładzie pracy / instytucji, opisać go, przedstawić koncepcję rozwiązania i ją zrealizować. & \\ \hline
09 & Potrafi rozwiązać rzeczywiste zadanie inżynierskie z zakresu działalności informatycznej zakładu pracy/instytucji stosując normy i standardy stosowane w informatyce oraz biorąc pod uwagę aspekty środowiskowe i etyczne. & \\ \hline
10 & Pracuje w zespole zajmującym się zawodowo branżą IT & \\ \hline
11 & Przestrzega zasad etyki zawodowej i zgodnie z tymi zasadami korzysta z wiedzy i pomocy doświadczonych kolegów & \\ \hline
12 & Kontaktując się z osobami spoza branży potrafi zarówno pozyskać od nich niezbędne informacje do realizacji planowanego zadania, jak i przekazać im w sposób zrozumiały informacje i opinie z zakresu informatyki & \\ \hline
13 & Dostrzega w praktyce tempo deaktualizacji wiedzy informatycznej oraz skutki działalności informatyków w szczególności ekonomiczne i społeczne & \\ \hline
\end{longtable}

\newpage

\begin{center}
\textbf{HARMONOGRAM PRAKTYKI ZAWODOWEJ}
\end{center}

\begin{center}
\renewcommand{\arraystretch}{1.5}
\begin{tabular}{|c|p{8cm}|p{4cm}|}
\hline
\textbf{L.p.} & \textbf{Dział / komórka (miejsce odbywania praktyki)} & \textbf{Planowana liczba dni roboczych} \\ \hline
1 & & \\ \hline
2 & & \\ \hline
3 & & \\ \hline
4 & & \\ \hline
5 & & \\ \hline
6 & & \\ \hline
7 & & \\ \hline
8 & & \\ \hline
9 & & \\ \hline
10 & & \\ \hline
11 & & \\ \hline
12 & & \\ \hline
13 & & \\ \hline
\multicolumn{2}{|r|}{\textbf{Łącznie}} & \\ \hline
\multicolumn{2}{|r|}{\textbf{Wymagana}} & \textbf{120} \\ \hline
\end{tabular}
\end{center}

\vspace{3em}
\noindent Uzgodniono w dniu: \dotfill

\vspace{4em}
\noindent
\begin{minipage}[t]{0.3\textwidth}
\centering
\dotfill\\
\small podpis uczelnianego opiekuna praktyki
\end{minipage}
\hfill
\begin{minipage}[t]{0.3\textwidth}
\centering
\dotfill\\
\small podpis zakładowego opiekuna praktyki
\end{minipage}
\hfill
\begin{minipage}[t]{0.3\textwidth}
\centering
\dotfill\\
\small podpis studenta
\end{minipage}

\end{document}